% Options for packages loaded elsewhere
% Options for packages loaded elsewhere
\PassOptionsToPackage{unicode}{hyperref}
\PassOptionsToPackage{hyphens}{url}
\PassOptionsToPackage{dvipsnames,svgnames,x11names}{xcolor}
%
\documentclass[
  11pt,
]{article}
\usepackage{xcolor}
\usepackage[margin=1in]{geometry}
\usepackage{amsmath,amssymb}
\setcounter{secnumdepth}{-\maxdimen} % remove section numbering
\usepackage{iftex}
\ifPDFTeX
  \usepackage[T1]{fontenc}
  \usepackage[utf8]{inputenc}
  \usepackage{textcomp} % provide euro and other symbols
\else % if luatex or xetex
  \usepackage{unicode-math} % this also loads fontspec
  \defaultfontfeatures{Scale=MatchLowercase}
  \defaultfontfeatures[\rmfamily]{Ligatures=TeX,Scale=1}
\fi
\usepackage{lmodern}
\ifPDFTeX\else
  % xetex/luatex font selection
\fi
% Use upquote if available, for straight quotes in verbatim environments
\IfFileExists{upquote.sty}{\usepackage{upquote}}{}
\IfFileExists{microtype.sty}{% use microtype if available
  \usepackage[]{microtype}
  \UseMicrotypeSet[protrusion]{basicmath} % disable protrusion for tt fonts
}{}
\makeatletter
\@ifundefined{KOMAClassName}{% if non-KOMA class
  \IfFileExists{parskip.sty}{%
    \usepackage{parskip}
  }{% else
    \setlength{\parindent}{0pt}
    \setlength{\parskip}{6pt plus 2pt minus 1pt}}
}{% if KOMA class
  \KOMAoptions{parskip=half}}
\makeatother
% Make \paragraph and \subparagraph free-standing
\makeatletter
\ifx\paragraph\undefined\else
  \let\oldparagraph\paragraph
  \renewcommand{\paragraph}{
    \@ifstar
      \xxxParagraphStar
      \xxxParagraphNoStar
  }
  \newcommand{\xxxParagraphStar}[1]{\oldparagraph*{#1}\mbox{}}
  \newcommand{\xxxParagraphNoStar}[1]{\oldparagraph{#1}\mbox{}}
\fi
\ifx\subparagraph\undefined\else
  \let\oldsubparagraph\subparagraph
  \renewcommand{\subparagraph}{
    \@ifstar
      \xxxSubParagraphStar
      \xxxSubParagraphNoStar
  }
  \newcommand{\xxxSubParagraphStar}[1]{\oldsubparagraph*{#1}\mbox{}}
  \newcommand{\xxxSubParagraphNoStar}[1]{\oldsubparagraph{#1}\mbox{}}
\fi
\makeatother


\usepackage{longtable,booktabs,array}
\usepackage{calc} % for calculating minipage widths
% Correct order of tables after \paragraph or \subparagraph
\usepackage{etoolbox}
\makeatletter
\patchcmd\longtable{\par}{\if@noskipsec\mbox{}\fi\par}{}{}
\makeatother
% Allow footnotes in longtable head/foot
\IfFileExists{footnotehyper.sty}{\usepackage{footnotehyper}}{\usepackage{footnote}}
\makesavenoteenv{longtable}
\usepackage{graphicx}
\makeatletter
\newsavebox\pandoc@box
\newcommand*\pandocbounded[1]{% scales image to fit in text height/width
  \sbox\pandoc@box{#1}%
  \Gscale@div\@tempa{\textheight}{\dimexpr\ht\pandoc@box+\dp\pandoc@box\relax}%
  \Gscale@div\@tempb{\linewidth}{\wd\pandoc@box}%
  \ifdim\@tempb\p@<\@tempa\p@\let\@tempa\@tempb\fi% select the smaller of both
  \ifdim\@tempa\p@<\p@\scalebox{\@tempa}{\usebox\pandoc@box}%
  \else\usebox{\pandoc@box}%
  \fi%
}
% Set default figure placement to htbp
\def\fps@figure{htbp}
\makeatother





\setlength{\emergencystretch}{3em} % prevent overfull lines

\providecommand{\tightlist}{%
  \setlength{\itemsep}{0pt}\setlength{\parskip}{0pt}}



 


\usepackage{amsmath,amssymb,enumitem,booktabs}
\setlist{itemsep=2pt,topsep=3pt}
\usepackage{fancyhdr}\pagestyle{fancy}\fancyhf{}
\fancyhead[L]{\small DS701 Fall 2026}\fancyhead[R]{\small Homework 2}
\fancyfoot[C]{\small\thepage}
\renewcommand{\headrulewidth}{0.4pt}
\usepackage[breakable]{tcolorbox}
% --- answer-space macros (problems PDF only; no-ops in the solutions PDF) ---
% \answerspace{2in}  vertical room for handwritten work, sized from the solution
% \blank             a short fill-in rule; \blank[1.5in] for a wider one
% \answerpage        start the next problem on a fresh page (Gradescope page matching)
\newcommand{\answerspace}[1]{\par\vspace*{#1}}  % \vspace* so room is not dropped at a page break
\newcommand{\blank}[1][1.2in]{\rule{#1}{0.4pt}}
\newcommand{\answerpage}{\clearpage}
\newtcolorbox{solutionbox}{breakable,colback=blue!4,colframe=blue!45!black,title=Solution,fonttitle=\bfseries,left=4pt,right=4pt,top=3pt,bottom=3pt}
\makeatletter
\@ifpackageloaded{caption}{}{\usepackage{caption}}
\AtBeginDocument{%
\ifdefined\contentsname
  \renewcommand*\contentsname{Table of contents}
\else
  \newcommand\contentsname{Table of contents}
\fi
\ifdefined\listfigurename
  \renewcommand*\listfigurename{List of Figures}
\else
  \newcommand\listfigurename{List of Figures}
\fi
\ifdefined\listtablename
  \renewcommand*\listtablename{List of Tables}
\else
  \newcommand\listtablename{List of Tables}
\fi
\ifdefined\figurename
  \renewcommand*\figurename{Figure}
\else
  \newcommand\figurename{Figure}
\fi
\ifdefined\tablename
  \renewcommand*\tablename{Table}
\else
  \newcommand\tablename{Table}
\fi
}
\@ifpackageloaded{float}{}{\usepackage{float}}
\floatstyle{ruled}
\@ifundefined{c@chapter}{\newfloat{codelisting}{h}{lop}}{\newfloat{codelisting}{h}{lop}[chapter]}
\floatname{codelisting}{Listing}
\newcommand*\listoflistings{\listof{codelisting}{List of Listings}}
\makeatother
\makeatletter
\makeatother
\makeatletter
\@ifpackageloaded{caption}{}{\usepackage{caption}}
\@ifpackageloaded{subcaption}{}{\usepackage{subcaption}}
\makeatother
\usepackage{bookmark}
\IfFileExists{xurl.sty}{\usepackage{xurl}}{} % add URL line breaks if available
\urlstyle{same}
\hypersetup{
  pdftitle={DS701 Homework 2: Distances},
  colorlinks=true,
  linkcolor={blue},
  filecolor={Maroon},
  citecolor={Blue},
  urlcolor={Blue},
  pdfcreator={LaTeX via pandoc}}


\title{DS701 Homework 2: Distances}
\usepackage{etoolbox}
\makeatletter
\providecommand{\subtitle}[1]{% add subtitle to \maketitle
  \apptocmd{\@title}{\par {\large #1 \par}}{}{}
}
\makeatother
\subtitle{Released Mon Sep 14 · due Mon Sep 21, 11:59 pm on Gradescope}
\author{}
\date{}
\begin{document}
\maketitle


\noindent\textbf{Name:}\ \rule{2.0in}{0.4pt}\qquad\textbf{BU ID:}\ \rule{1.2in}{0.4pt}
\vspace{6pt}

\textbf{How this works.} Answer \textbf{by hand on paper} (or a tablet
with a stylus), show your work, then scan and upload to Gradescope,
matching each problem to its pages. This is deliberate practice for the
closed-notes quizzes, which are in the same style --- so do it without a
computer or an AI assistant. It is \textbf{participation-graded}: a
serious attempt at every problem earns full credit, and worked solutions
are released the day after the deadline. Budget 1--2 hours.

\textbf{Covers:} Lecture 02 (Distances and Similarity --- metrics,
\(\ell_p\) distances, cosine, Hamming/Jaccard, feature scaling). Time
series and DTW are \emph{not} covered. Square roots may be left as
\(\sqrt{\cdot}\) with a one-decimal approximation next to them
(e.g.~\(\sqrt{12}\approx 3.5\)).

\subsection{Problem 1: The metric decides who is
close}\label{problem-1-the-metric-decides-who-is-close}

\emph{Practices: Minkowski distances, cosine similarity, Hamming and
Jaccard --- and noticing that ``closest'' depends on the choice.}

Three short documents are represented as word-count vectors over the
vocabulary (\emph{data}, \emph{model}, \emph{error}):
\[\mathbf{A} = (1, 1, 1),\qquad \mathbf{B} = (3, 3, 3),\qquad \mathbf{C} = (6, 1, 1).\]

\begin{enumerate}
\def\labelenumi{(\alph{enumi})}
\tightlist
\item
  For each of the three pairs (A,B), (A,C), (B,C) compute the Minkowski
  distance \(\Vert\mathbf{x}-\mathbf{y}\Vert_p\) for \(p = 1\)
  (Manhattan), \(p = 2\) (Euclidean) and \(p = \infty\) (Chebyshev) ---
  nine numbers, arranged as a \(3\times3\) table. Under each of the
  three distances, which pair is \textbf{closest}? Note where the answer
  changes.
\end{enumerate}

\begin{longtable}[]{@{}cccc@{}}
\toprule\noalign{}
pair & \(p=1\) & \(p=2\) & \(p=\infty\) \\
\midrule\noalign{}
\endhead
\bottomrule\noalign{}
\endlastfoot
(A,B) & \blank[0.8in] & \blank[0.8in] & \blank[0.8in] \\
(A,C) & \blank[0.8in] & \blank[0.8in] & \blank[0.8in] \\
(B,C) & \blank[0.8in] & \blank[0.8in] & \blank[0.8in] \\
\end{longtable}

Closest pair: \(p=1\): \blank[0.7in] \(\quad p=2\):
\blank[0.7in] \(\quad p=\infty\): \blank[0.7in]

\answerspace{2.2in}

\begin{enumerate}
\def\labelenumi{(\alph{enumi})}
\setcounter{enumi}{1}
\tightlist
\item
  Compute the cosine similarity for the same three pairs (three
  decimals; \(\sqrt{3}\approx1.732\), \(\sqrt{27}\approx5.196\),
  \(\sqrt{38}\approx6.164\), \(\sqrt{114}\approx10.677\)). Which pair is
  most similar under cosine? Explain in one sentence why cosine gives A
  and B the value it does even though every \(\ell_p\) distance in (a)
  says they are some way apart, and say which of the two views is the
  right one if we want to know whether two documents are \emph{about the
  same thing}.
\end{enumerate}

\(\cos(\mathbf{A},\mathbf{B}) =\)
\blank[0.8in] \(\qquad \cos(\mathbf{A},\mathbf{C}) =\)
\blank[0.8in] \(\qquad \cos(\mathbf{B},\mathbf{C}) =\) \blank[0.8in]

Most similar pair: \blank[0.8in]

\answerspace{2.4in}

\newpage % (c) with its table and answer room starts on a fresh page rather than splitting

\begin{enumerate}
\def\labelenumi{(\alph{enumi})}
\setcounter{enumi}{2}
\tightlist
\item
  Two pairs of shopping baskets, as sets of items:
\end{enumerate}

\begin{itemize}
\tightlist
\item
  \(P = \{\)bread, milk, eggs, butter, cheese, apples, rice,
  coffee\(\}\) and \(Q = \{\)bread, milk, eggs, butter, cheese, apples,
  rice, tea\(\}\);
\item
  \(R = \{\)bread, milk\(\}\) and \(S = \{\)bread, jam\(\}\).
\end{itemize}

For each pair compute the Hamming distance (the number of items in one
basket but not the other, i.e.~the size of the symmetric difference) and
the Jaccard similarity \(J_{\rm Sim} = |x\cap y|/|x\cup y|\). Which pair
does Hamming say is more alike, and which does Jaccard say? In one or
two sentences, say which measure you would trust for ``are these two
shoppers similar?'' and why.

\begin{longtable}[]{@{}ccc@{}}
\toprule\noalign{}
pair & Hamming & \(J_{\rm Sim}\) \\
\midrule\noalign{}
\endhead
\bottomrule\noalign{}
\endlastfoot
(P,Q) & \blank[0.8in] & \blank[0.8in] \\
(R,S) & \blank[0.8in] & \blank[0.8in] \\
\end{longtable}

\answerspace{1.8in}

\answerpage

\subsection{Problem 2: Is it a metric?}\label{problem-2-is-it-a-metric}

\emph{Practices: the metric axioms, especially the triangle inequality,
and why they matter downstream.}

Recall the four metric axioms: \(d(x,x)=0\); \(d(x,y)>0\) for
\(x\ne y\); \(d(x,y)=d(y,x)\); and the triangle inequality
\(d(x,y)\le d(x,z)+d(z,y)\) for all \(x,y,z\).

\begin{enumerate}
\def\labelenumi{(\alph{enumi})}
\tightlist
\item
  Let \(P=(0,0)\), \(Q=(3,4)\), \(R=(6,0)\) in \(\mathbb{R}^2\). Compute
  the three Euclidean distances \(d(P,Q)\), \(d(Q,R)\), \(d(P,R)\) and
  verify the triangle inequality for this triple in all three ways (each
  distance is at most the sum of the other two). Then state, in one
  line, when the triangle inequality holds with \textbf{equality} for
  Euclidean distance, and give a triple of points that achieves it.
\end{enumerate}

\(d(P,Q) =\) \blank[0.7in] \(\quad d(Q,R) =\)
\blank[0.7in] \(\quad d(P,R) =\) \blank[0.7in]

\answerspace{1.9in}

\begin{enumerate}
\def\labelenumi{(\alph{enumi})}
\setcounter{enumi}{1}
\tightlist
\item
  Cosine \emph{distance} is
  \(d_{\cos}(\mathbf{x},\mathbf{y}) = 1 - \cos(\mathbf{x},\mathbf{y})\).
  The lecture says it is \textbf{not} a metric. Prove it: give three
  explicit vectors \(\mathbf{x},\mathbf{z},\mathbf{y}\) in
  \(\mathbb{R}^2\) for which
  \(d_{\cos}(\mathbf{x},\mathbf{y}) > d_{\cos}(\mathbf{x},\mathbf{z}) + d_{\cos}(\mathbf{z},\mathbf{y})\),
  and show the numbers (use \(\cos 45^\circ = 1/\sqrt2 \approx 0.707\)
  if it helps). \emph{Hint:} think about angles \(0^\circ\),
  \(45^\circ\), \(90^\circ\).
\end{enumerate}

\(\mathbf{x} =\) \blank[0.9in] \(\quad \mathbf{z} =\)
\blank[0.9in] \(\quad \mathbf{y} =\) \blank[0.9in]

\answerspace{2.3in}

\begin{enumerate}
\def\labelenumi{(\alph{enumi})}
\setcounter{enumi}{2}
\tightlist
\item
  Cosine distance fails a \textbf{second} axiom as well. Which one, and
  what is a two-vector example? (One line.)
\end{enumerate}

\answerspace{1in}

\begin{enumerate}
\def\labelenumi{(\alph{enumi})}
\setcounter{enumi}{3}
\tightlist
\item
  In two or three sentences: why does it matter, for nearest-neighbor
  search or for clustering, whether the dissimilarity you use satisfies
  the triangle inequality? Concretely, what can you conclude about
  \(d(q,b)\) from knowing \(d(q,a)\) is small and \(d(a,b)\) is large
  when \(d\) is a metric --- and what goes wrong when it is not?
\end{enumerate}

\answerspace{1.8in}

\answerpage

\subsection{Problem 3: Same points, different units, different
neighbor}\label{problem-3-same-points-different-units-different-neighbor}

\emph{Practices: why raw Euclidean distance is dominated by the
largest-range feature, and what standardizing does.}

A bank describes customers by two features: annual income (in dollars)
and age (in years). A query customer \(q\) and four customers on file:

\begin{longtable}[]{@{}crr@{}}
\toprule\noalign{}
& income & age \\
\midrule\noalign{}
\endhead
\bottomrule\noalign{}
\endlastfoot
\(q\) & 50,000 & 30 \\
\(P_1\) & 52,000 & 60 \\
\(P_2\) & 60,000 & 32 \\
\(P_3\) & 45,000 & 28 \\
\(P_4\) & 51,000 & 45 \\
\end{longtable}

To \textbf{standardize} a feature, replace each value \(v\) by its
\(z\)-score \((v-\mu)/\sigma\), where \(\mu\) and \(\sigma\) are the
feature's mean and standard deviation. Use these values, taken from the
bank's full customer base: income \(\mu = 50{,}000\),
\(\sigma = 10{,}000\); age \(\mu = 40\), \(\sigma = 10\).

\begin{enumerate}
\def\labelenumi{(\alph{enumi})}
\tightlist
\item
  Using raw Euclidean distance, which of \(P_1,\dots,P_4\) is the
  nearest neighbor of \(q\)? (Squared distances are enough to rank;
  report them.) In one line: how much does the age difference contribute
  to any of these distances?
\end{enumerate}

\begin{longtable}[]{@{}
  >{\centering\arraybackslash}p{(\linewidth - 8\tabcolsep) * \real{0.2000}}
  >{\centering\arraybackslash}p{(\linewidth - 8\tabcolsep) * \real{0.2000}}
  >{\centering\arraybackslash}p{(\linewidth - 8\tabcolsep) * \real{0.2000}}
  >{\centering\arraybackslash}p{(\linewidth - 8\tabcolsep) * \real{0.2000}}
  >{\centering\arraybackslash}p{(\linewidth - 8\tabcolsep) * \real{0.2000}}@{}}
\toprule\noalign{}
\begin{minipage}[b]{\linewidth}\centering
\end{minipage} & \begin{minipage}[b]{\linewidth}\centering
\(P_1\)
\end{minipage} & \begin{minipage}[b]{\linewidth}\centering
\(P_2\)
\end{minipage} & \begin{minipage}[b]{\linewidth}\centering
\(P_3\)
\end{minipage} & \begin{minipage}[b]{\linewidth}\centering
\(P_4\)
\end{minipage} \\
\midrule\noalign{}
\endhead
\bottomrule\noalign{}
\endlastfoot
\(d^2\) from \(q\) & \blank[0.9in] & \blank[0.9in] & \blank[0.9in] &
\blank[0.9in] \\
\end{longtable}

Nearest neighbor: \blank[0.7in]

\answerspace{1.4in}

\begin{enumerate}
\def\labelenumi{(\alph{enumi})}
\setcounter{enumi}{1}
\tightlist
\item
  Standardize all five rows with the given \(\mu,\sigma\) (write the
  \(z\)-scores), recompute the four squared Euclidean distances from
  \(q\), and give the nearest neighbor now. Does it change?
\end{enumerate}

\begin{longtable}[]{@{}
  >{\centering\arraybackslash}p{(\linewidth - 10\tabcolsep) * \real{0.1667}}
  >{\centering\arraybackslash}p{(\linewidth - 10\tabcolsep) * \real{0.1667}}
  >{\centering\arraybackslash}p{(\linewidth - 10\tabcolsep) * \real{0.1667}}
  >{\centering\arraybackslash}p{(\linewidth - 10\tabcolsep) * \real{0.1667}}
  >{\centering\arraybackslash}p{(\linewidth - 10\tabcolsep) * \real{0.1667}}
  >{\centering\arraybackslash}p{(\linewidth - 10\tabcolsep) * \real{0.1667}}@{}}
\toprule\noalign{}
\begin{minipage}[b]{\linewidth}\centering
\end{minipage} & \begin{minipage}[b]{\linewidth}\centering
\(q\)
\end{minipage} & \begin{minipage}[b]{\linewidth}\centering
\(P_1\)
\end{minipage} & \begin{minipage}[b]{\linewidth}\centering
\(P_2\)
\end{minipage} & \begin{minipage}[b]{\linewidth}\centering
\(P_3\)
\end{minipage} & \begin{minipage}[b]{\linewidth}\centering
\(P_4\)
\end{minipage} \\
\midrule\noalign{}
\endhead
\bottomrule\noalign{}
\endlastfoot
\(z_{\text{income}}\) & \blank[0.7in] & \blank[0.7in] & \blank[0.7in] &
\blank[0.7in] & \blank[0.7in] \\
\(z_{\text{age}}\) & \blank[0.7in] & \blank[0.7in] & \blank[0.7in] &
\blank[0.7in] & \blank[0.7in] \\
\(d^2\) from \(q\) & --- & \blank[0.7in] & \blank[0.7in] & \blank[0.7in]
& \blank[0.7in] \\
\end{longtable}

Nearest neighbor: \blank[0.7in]

\answerspace{2in}

\begin{enumerate}
\def\labelenumi{(\alph{enumi})}
\setcounter{enumi}{2}
\tightlist
\item
  Redo (a) with income measured in \emph{thousands} of dollars (50, 52,
  60, 45, 51) and age unchanged. Which neighbor wins? Then answer: which
  of the three answers should a data scientist trust as ``the customer
  most like \(q\)'', and why --- what is wrong with the other two?
\end{enumerate}

\begin{longtable}[]{@{}
  >{\centering\arraybackslash}p{(\linewidth - 8\tabcolsep) * \real{0.2000}}
  >{\centering\arraybackslash}p{(\linewidth - 8\tabcolsep) * \real{0.2000}}
  >{\centering\arraybackslash}p{(\linewidth - 8\tabcolsep) * \real{0.2000}}
  >{\centering\arraybackslash}p{(\linewidth - 8\tabcolsep) * \real{0.2000}}
  >{\centering\arraybackslash}p{(\linewidth - 8\tabcolsep) * \real{0.2000}}@{}}
\toprule\noalign{}
\begin{minipage}[b]{\linewidth}\centering
\end{minipage} & \begin{minipage}[b]{\linewidth}\centering
\(P_1\)
\end{minipage} & \begin{minipage}[b]{\linewidth}\centering
\(P_2\)
\end{minipage} & \begin{minipage}[b]{\linewidth}\centering
\(P_3\)
\end{minipage} & \begin{minipage}[b]{\linewidth}\centering
\(P_4\)
\end{minipage} \\
\midrule\noalign{}
\endhead
\bottomrule\noalign{}
\endlastfoot
\(d^2\) from \(q\) & \blank[0.9in] & \blank[0.9in] & \blank[0.9in] &
\blank[0.9in] \\
\end{longtable}

Nearest neighbor: \blank[0.7in] \(\quad\) Trust: (a) / (b) / (c)

\answerspace{2in}

\begin{enumerate}
\def\labelenumi{(\alph{enumi})}
\setcounter{enumi}{3}
\tightlist
\item
  Standardizing is not always right. Give one situation where you would
  \textbf{not} standardize before computing distances, and one sentence
  on why (think about features that are already in the same meaningful
  units, or a feature whose spread \emph{should} count).
\end{enumerate}

\answerspace{1.8in}




\end{document}
